\documentclass[11pt]{article}
\usepackage[margin=1in]{geometry}
\usepackage{amsmath,amssymb,amsthm}
\usepackage[hidelinks]{hyperref}
\newtheorem{theorem}{Theorem}
\newtheorem{lemma}[theorem]{Lemma}
\newtheorem{corollary}[theorem]{Corollary}
\theoremstyle{remark}\newtheorem{remark}[theorem]{Remark}
\newcommand{\F}{\mathbb F}
\newcommand{\topc}{\operatorname{top}}
\newcommand{\red}{\operatorname{red}}
\title{Maximum-degree local permutation polynomials over the prime field}
\author{}
\date{September 8, 2026}
\begin{document}
\maketitle
\begin{abstract}
For every prime power $q=p^r>3$ and every positive integer $n$, we prove the
existence of a local permutation polynomial over $\F_q$ with coefficients in
$\F_p$ and total degree $n(q-2)$. The proof uses a subfield patch that preserves
the extremal coefficient, together with an analysis of a recursive family
formed from inversion and a transposition. For every nonprime odd field, we
compute the extremal coefficient of that family in every arity $n\geq2$ by a first-order
recurrence over $\F_p$. In particular, over every $\F_{5^r}$ with $r\geq2$, the
family fails to have maximum degree among arities $n\geq2$ precisely when
$n$ is congruent to $6$ or $7$ modulo $10$. This gives infinite families of counterexamples to the proposed
maximum-degree conjecture for this recursion, while the ternary case supplies
the missing construction for the uniform existence result.
\end{abstract}

\section{Introduction and main results}

A \emph{local permutation polynomial} (LPP) over a finite field $E=\F_q$ is a
polynomial function $F:E^n\to E$ that is a permutation in each variable when the
other variables are fixed. We identify polynomial functions with their unique
reduced representatives, whose degree in each variable is less than $q$.
For $q>2$, an LPP has degree at most $q-2$ in each variable, and hence total
degree at most $n(q-2)$. We call this total degree \emph{maximum degree}.

Guti\'errez and Jim\'enez Urroz established unrestricted-coefficient existence
and several prime-coefficient cases in \cite{GJU-afr,GJU-mal}. Their remaining
prime-coefficient existence question is answered by the following theorem.

\begin{theorem}\label{thm:exist}
For every prime power $q=p^r>3$ and every integer $n\geq1$, there is an LPP in
$\F_p[X_1,\ldots,X_n]$ over $\F_q$ of reduced total degree $n(q-2)$.
\end{theorem}

The characteristic-two case is already known \cite[Theorem 4]{GJU-afr}.
For $p\geq5$, we extend prime-field examples by a patching construction that
preserves the extremal coefficient. The case $p=3$ follows from the next
result, which also describes the failure of a proposed general construction.

Let $q>3$, put $I(z)=z^{q-2}$, and let
\[
 t(z)=z+\sum_{j=0}^{q-2}z^j.
\]
As functions, $I$ inverts nonzero elements and fixes zero, while $t$ interchanges
zero and one and fixes all other elements. Define reduced polynomial functions
\begin{equation}\label{eq:family}
 f_1=X_1,\qquad
 f_n=t\bigl(I(f_{n-1})+I(X_n)\bigr)\quad(n\geq2).
\end{equation}
Every $f_n$ is an LPP defined over $\F_p$. The assertion that this recursion
always has maximum degree in odd characteristic is Conjecture 1 of both
\cite[Section 6]{GJU-afr} and \cite[Section 4]{GJU-mal}.

For a reduced polynomial $F$ in $n$ variables, write
\[
 \topc_q(F)=[X_1^{q-2}\cdots X_n^{q-2}]F.
\]
For an LPP, maximum degree is equivalent to $\topc_q(F)\ne0$.

\begin{theorem}\label{thm:recurrence}
Let $q=p^r$ with $p$ odd and $r\geq2$. Set $C_n=\topc_q(f_n)$ for the family
\eqref{eq:family}. Define $E_k\in\F_p$ by
\begin{equation}\label{eq:Ek}
 E_1=2,\qquad E_{k+1}=-4E_k+4.
\end{equation}
Then, for every $k\geq1$,
\begin{equation}\label{eq:Cn}
 C_{2k}=E_k,\qquad C_{2k+1}=-2E_k.
\end{equation}
In particular, $f_{2k}$ and $f_{2k+1}$ have maximum degree if and only if
$E_k\ne0$. Explicitly,
\begin{equation}\label{eq:closed}
 E_k=\begin{cases}
 \bigl(6(-4)^{k-1}+4\bigr)/5,&p\ne5,\\
 4k-2,&p=5,
 \end{cases}
\end{equation}
where all quantities are interpreted in $\F_p$.
\end{theorem}

\begin{corollary}\label{cor:failure}
For the recursion \eqref{eq:family} over $\F_{p^r}$, $r\geq2$, the following hold.
\begin{enumerate}
\item If $p=3$, then $C_n=2$ for every $n\geq2$.
\item If $p=5$, then $f_n$ fails to have maximum degree precisely when
$n\equiv6$ or $7\pmod {10}$, for $n\geq2$.
\item If $p\geq7$, then the failures in the paired arities $2k,2k+1$ occur
precisely when $(-4)^{k-1}=-2/3$ in $\F_p$. Either there are no failures or the
corresponding $k$ form one residue class modulo the multiplicative order of
$-4$ in $\F_p^\times$.
\end{enumerate}
\end{corollary}

In particular, $q=25$ and $n=6$ give a counterexample to the recursive
conjecture with degree strictly smaller than $138$. Theorem~\ref{thm:recurrence}
classifies maximum-degree arities, rather than the precise degree when the
extremal coefficient vanishes. Its hypothesis $r\geq2$ is essential to the
coefficient-counting argument below; no classification over prime fields is
asserted here. Also, $f_1=X_1$ is not of maximum degree for $q>3$, so our
statements about the recursion start at arity two.

\section{Extremal coefficients and subfield patching}

We first record the coefficient functional used throughout the paper.
In two variables, replacing an embedded subsquare is a form of Latin trade;
see \cite[Section 1]{CF} for the trade viewpoint. The feature needed here is
the exact preservation of a polynomial coefficient across different fields,
together with Frobenius equivariance.

\begin{lemma}\label{lem:coefficient}
Let $q>2$. For every reduced polynomial function $P:\F_q^n\to\F_q$,
\begin{equation}\label{eq:moment}
 \topc_q(P)=(-1)^n\sum_{a\in\F_q^n}
 \left(\prod_{i=1}^n a_i\right)P(a).
\end{equation}
If $P$ is an LPP, then its degree in each variable is at most $q-2$.
\end{lemma}

\begin{proof}
For a reduced univariate polynomial $h$, the sums of powers in $\F_q$ give
\[
 \sum_{a\in\F_q}a h(a)=-[X^{q-2}]h.
\]
Indeed, for $0\leq j\leq q-1$, the sum of $a^{j+1}$ is nonzero only for
$j=q-2$, when it equals $-1$. Applying this identity one variable at a time
proves \eqref{eq:moment}.

On a coordinate line, an LPP takes each field value once. Its sum of values is
therefore zero, as $q>2$. For a reduced univariate polynomial this sum equals
minus the coefficient of $X^{q-1}$. Consequently the coefficient polynomial
of $X_i^{q-1}$ in $P$ vanishes on every choice of the remaining coordinates.
It is reduced in those coordinates and hence is the zero polynomial.
\end{proof}

\begin{lemma}[Subfield patch]\label{lem:patch}
Let $k=\F_s\subseteq E=\F_q$ with $s>3$, and let $h:k^n\to k$ be an LPP of
maximum degree. Define
\begin{equation}\label{eq:patch}
 F(x_1,\ldots,x_n)=\begin{cases}
 h(x_1,\ldots,x_n),&(x_1,\ldots,x_n)\in k^n,\\
 x_1+\cdots+x_n,&(x_1,\ldots,x_n)\notin k^n.
 \end{cases}
\end{equation}
Then $F$ is an LPP of maximum degree over $E$, and
\begin{equation}\label{eq:patchcoefficient}
 \topc_q(F)=\topc_s(h).
\end{equation}
If $h$ is equivariant under the prime-field Frobenius automorphism, then so is
$F$, and its reduced polynomial has coefficients in the prime field.
\end{lemma}

\begin{proof}
Fix all coordinates except $x_i$. If a fixed coordinate lies outside $k$, the
coordinate map is translation by a constant of $E$. Otherwise all fixed
coordinates lie in $k$. On $k$ the coordinate map is a permutation of $k$,
because $h$ is an LPP. On $E\setminus k$ it is translation by a constant in
$k$, which permutes $E\setminus k$. These two parts have disjoint images, so
the coordinate map permutes $E$.

Put $A(x)=x_1+\cdots+x_n$. Since $s>3$ and $q\geq s$, both $\topc_s(A)$ and
$\topc_q(A)$ vanish: the relevant total degree is greater than one. The
function $F-A$ is supported on $k^n$. Lemma~\ref{lem:coefficient} therefore
gives
\[
 \topc_q(F)=(-1)^n\sum_{a\in k^n}\left(\prod_i a_i\right)(h(a)-A(a))
 =\topc_s(h)\ne0.
\]
The degree bound in that lemma proves maximality.

Finally $k$ is stable under the prime-field Frobenius, and addition commutes
with Frobenius. Thus \eqref{eq:patch} is equivariant whenever $h$ is. To see
that equivariance is equivalent to prime-field coefficients, write the reduced
representative as $P=\sum_\alpha c_\alpha X^\alpha$. Its coefficientwise
Frobenius conjugate is $P^\sigma=\sum_\alpha c_\alpha^pX^\alpha$.
Equivariance and the bijectivity of Frobenius imply that $P$ and $P^\sigma$
agree at every point of $E^n$. Uniqueness of reduced interpolation yields
$c_\alpha^p=c_\alpha$ for all $\alpha$.
\end{proof}

The patch has a direct polynomial representative. If $h$ denotes its reduced
polynomial over $k$, then
\begin{equation}\label{eq:explicitpatch}
 F=\red_q\left(
 A+\prod_{i=1}^n\bigl(1-(X_i^s-X_i)^{q-1}\bigr)(h-A)
 \right).
\end{equation}
Here $\red_q$ means reduction modulo $(X_1^q-X_1,\ldots,X_n^q-X_n)$.
Each bracket is the indicator function of $k$. In particular, when $s=p$,
formula \eqref{eq:explicitpatch} visibly has coefficients in $\F_p$.

\section{A transfer operator for the recursion}

Let $E=\F_q$ with $q>3$. On the space of reduced univariate polynomial
functions $h:E\to E$, define the linear operator
\begin{equation}\label{eq:T}
 (Th)(z)=\sum_{x\in E}x\,h\bigl(t(I(z)+I(x))\bigr).
\end{equation}
For $n\geq1$ put
\[
 S_n(h)=\sum_{x\in E^n}\left(\prod_{i=1}^n x_i\right)h(f_n(x)).
\]
Summing in the last coordinate of \eqref{eq:family} shows that
$S_n(h)=S_{n-1}(Th)$. Since $S_1(h)=\sum_z zh(z)=-[X^{q-2}]h$, we obtain
\begin{equation}\label{eq:transfercoefficient}
 C_n=(-1)^{n+1}[X^{q-2}]T^{n-1}X.
\end{equation}

\begin{lemma}\label{lem:T}
Let $h=\sum_{j=0}^{q-1}h_jX^j$, and put
$\Delta=h(1)-h(0)=\sum_{j=1}^{q-1}h_j$. The coefficients of $Th$ satisfy
\begin{align}
 (Th)_0&=-h_1-\Delta, & (Th)_1&=h_{q-1},\label{eq:Tend}\\
 (Th)_i&=i(h_{q-i}+\Delta)\quad(2\leq i\leq q-2),
 & (Th)_{q-1}&=0.\label{eq:Tmiddle}
\end{align}
The integer $i$ in \eqref{eq:Tmiddle} is interpreted as an element of $E$.
\end{lemma}

\begin{proof}
Let $g$ be the reduced representative of $h\circ t$. Since $t$ swaps zero and
one, the difference between $g$ and $h$ is represented by
\[
 g(X)=h(X)+\Delta\sum_{j=0}^{q-2}X^j.
\]
Indeed the sum on the right is $1$ at zero, $-1$ at one, and zero elsewhere.
In \eqref{eq:T}, substitute $u=I(x)$ and write $v=I(z)$. Expansion of $g(v+u)$
gives
\[
 (Th)(z)=\sum_{u\in E}I(u)g(v+u)=-g'(v).
\]
For clarity, the coefficient of $u^j$ in that expansion is multiplied by
$\sum_{u\ne0}u^{j-1}$. Among $0\leq j\leq q-1$, this sum vanishes except
at $j=1$, where it is $-1$. The surviving coefficient is precisely $g'(v)$.

Thus $Th=-g'\circ I$. Its constant coefficient is $-g_1=-h_1-\Delta$.
For $2\leq j\leq q-1$, the monomial
$-j g_jI(X)^{j-1}$ reduces to $-j g_jX^{q-j}$.
When $j=q-1$ this gives coefficient $h_{q-1}$ at $X$; the other cases give
\eqref{eq:Tmiddle}. No term of degree $q-1$ occurs.
\end{proof}

\section{The extremal coefficient over nonprime odd fields}

We now prove Theorem~\ref{thm:recurrence}. Assume $q=p^r$ with $p$ odd and
$r\geq2$, and put $H_j=T^jX$ for $j\geq1$.
Lemma~\ref{lem:T} shows that the coefficients of $H_j$ at indices $1$ and
$q-1$ are zero. Its remaining positive-degree coefficients have the form
\begin{equation}\label{eq:profile}
 [X^i]H_j=P_j(i\bmod p)\quad(2\leq i\leq q-2),
\end{equation}
where $P_j:\F_p\to\F_p$ and $P_j(0)=0$. For $j=1$ we may take $P_1(a)=a$;
the profile property then follows inductively from \eqref{eq:Tmiddle}.

Write $\Delta_j=H_j(1)-H_j(0)$. In the complete index interval
$0,\ldots,q-1$, each residue modulo $p$ occurs $p^{r-1}$ times. This multiplicity
is zero in $\F_p$. Removing indices $0,1,q-1$ and using $P_j(0)=0$ gives
\begin{equation}\label{eq:Deltaj}
 \Delta_j=-P_j(1)-P_j(-1).
\end{equation}
By Lemma~\ref{lem:T} we can therefore choose the profiles recursively as
\begin{equation}\label{eq:Pupdate}
 P_{j+1}(a)=a\bigl(P_j(-a)+\Delta_j\bigr).
\end{equation}
Evaluating at $a=1$ and $a=-1$ yields
\[
 P_{j+1}(1)=-P_j(1),\qquad P_{j+1}(-1)=P_j(-1).
\]
Since $P_1(1)=1$ and $P_1(-1)=-1$, we conclude that
\begin{equation}\label{eq:Deltaalternates}
 P_j(1)=(-1)^{j-1},\qquad P_j(-1)=-1,\qquad
 \Delta_j=1-(-1)^{j-1}.
\end{equation}

Put $a_j=P_j(-2)$ and $b_j=P_j(2)$. Equation~\eqref{eq:Pupdate} gives
\[
 a_{j+1}=-2(b_j+\Delta_j),\qquad
 b_{j+1}=2(a_j+\Delta_j).
\]
Combining the two identities and substituting \eqref{eq:Deltaalternates},
\begin{equation}\label{eq:aj}
 a_{j+2}=-4a_j-4\Delta_j-2\Delta_{j+1}
 =-4a_j-6+2(-1)^{j-1}.
\end{equation}
These evaluations remain valid when $p=3$, where some of the named residues
coincide: they are values of the same functions, not coordinates required to
be distinct.

The residue of $q-2$ modulo $p$ is $-2$, so
\eqref{eq:transfercoefficient} becomes $C_n=(-1)^{n+1}a_{n-1}$ for $n\geq2$.
Consequently \eqref{eq:aj} gives
\begin{equation}\label{eq:Crec}
 C_{n+2}=-4C_n+6(-1)^n-2\quad(n\geq2).
\end{equation}
The initial values are $C_2=2$ and $C_3=-4$, obtained from $P_1(a)=a$ and
$P_2(a)=-a^2$. Thus the even subsequence is the sequence $E_k$ in
\eqref{eq:Ek}. The odd subsequence starts at $-2E_1$ and obeys
$C_{2k+3}=-4C_{2k+1}-8$, which proves $C_{2k+1}=-2E_k$ by induction.
Solving \eqref{eq:Ek} proves \eqref{eq:closed}. Lemma~\ref{lem:coefficient}
then proves the maximum-degree criterion, completing
Theorem~\ref{thm:recurrence}. Corollary~\ref{cor:failure} follows immediately.

\section{Uniform existence over the prime field}

We complete the proof of Theorem~\ref{thm:exist} and give formulas for its
constructions. For $n=1$, the polynomial $X^{q-2}$ works in every field under
consideration. It remains to treat $n\geq2$.

If $p=2$, the known construction \cite[Theorem 4]{GJU-afr} is
\begin{equation}\label{eq:binary}
 B_{q,n}(X_1,\ldots,X_n)
 =\prod_{i=1}^n\left(\sum_{j=1}^{q-2}X_i^j\right)+\sum_{i=1}^nX_i.
\end{equation}
It is an LPP of degree $n(q-2)$, with coefficients in $\F_2$.

Suppose $p\geq5$. Theorem 5 of \cite{GJU-afr}, applied over $\F_p$ with
$b=p-2$, provides a maximum-degree LPP $h$ in every arity: the required
conditions are $1<b<p-1$ and $\gcd(b,p-1)=1$. For $q=p$ this is the desired
polynomial. For any larger extension, apply Lemma~\ref{lem:patch} with
$s=p$. Formula~\eqref{eq:explicitpatch} gives the resulting polynomial over
$\F_p$ and \eqref{eq:patchcoefficient} proves its maximum degree.

For completeness, the prime-field seed can be chosen by a finite explicit
procedure from that construction. Set $G_0(Y)=Y$. Form $G_{d+1}$ by taking
$b$ disjoint copies of $G_d$, adding them, and raising the sum to the $b$th
power. Choose $d$ with $N=b^d\geq n$, and form
\[
 H=\red_p G_d(Y_1^{p-2},\ldots,Y_N^{p-2}).
\]
The seed construction proves that $H$ has nonzero coefficient at
$\prod_{i=1}^N Y_i^{p-2}$. Specialize the last $N-n$ variables to a tuple
in $\F_p^{N-n}$ for which the coefficient of
$\prod_{i=1}^nY_i^{p-2}$ remains nonzero. Such a tuple exists because that
coefficient is a nonzero reduced polynomial in the last variables. Choosing
the lexicographically first such tuple, with field elements represented by
$0,1,\ldots,p-1$, makes this a deterministic procedure. Specialization
preserves the local permutation property.

Finally, suppose $p=3$. Since $q>3$, its extension degree satisfies $r\geq2$.
Use the recursion \eqref{eq:family}. Corollary~\ref{cor:failure} shows that
$\topc_q(f_n)=2\ne0$ for every $n\geq2$, and every operation in the recursion
has coefficients in $\F_3$. This proves Theorem~\ref{thm:exist}.

\section{Scope and further questions}

The subfield patch separates coefficient descent from the original
construction of an LPP. Its extremal coefficient is preserved exactly, rather
than merely shown to remain nonzero. The ternary recursion supplies the case
where the prime subfield itself is too small for the patching lemma.

The transfer identity $Th=-(h\circ t)'\circ I$, where the derivative is taken
after reducing $h\circ t$, also applies when $q$ is prime. The simplification
\eqref{eq:Deltaj} does not: residue classes then occur only once. Determining
the maximum-degree arities in that case is a separate problem. For nonprime
odd fields, another question is to determine the exact reduced degree at
the arities where Theorem~\ref{thm:recurrence} gives $C_n=0$.

\begin{thebibliography}{9}
\bibitem{CF}
N.~J.~Cavenagh and R.~M.~Falc\'on,
\emph{Latin bitrades derived from quasigroup autoparatopisms},
arXiv:2308.14987v1 (2023).
\url{https://arxiv.org/abs/2308.14987}.

\bibitem{GJU-afr}
J.~Gutierrez and J.~Jim\'enez Urroz,
\emph{Permutation and local permutation polynomials of maximum degree},
Afrika Matematika \textbf{36} (2025), article 45.
\href{https://doi.org/10.1007/s13370-025-01247-3}{doi:10.1007/s13370-025-01247-3}.

\bibitem{GJU-mal}
J.~Gutierrez and J.~Jim\'enez Urroz,
\emph{Local permutation polynomials of maximum degree over prime finite fields},
Bulletin of the Malaysian Mathematical Sciences Society \textbf{48} (2025),
article 40.
\href{https://doi.org/10.1007/s40840-025-01825-5}{doi:10.1007/s40840-025-01825-5}.
\end{thebibliography}
\end{document}
